%%%%%%%%%%  scan idx 100  |  folio 89  %%%%%%%%%%

\underline{Corolario 3.13.}

\begin{enumerate}[label=(\alph*)]
\item El $R$-álgebra $A$ es formalmente lisa sobre $R$ si y solo si
$D^1(A|R,M)=0$ para todo $A$-módulo $M$.

\item Sea $R$ noetheriano y sea $A$ un $R$-álgebra de tipo finito.
Entonces $A$ es una intersección completa relativa sobre $R$ si y solo
si $D^2(A|R,M)=0$ para todo $A$-módulo $M$.

\item Sea $R$ noetheriano y sea $A$ un $R$-álgebra de tipo finito.
Para todo subconjunto multiplicativo $S$ de $A$, hay un isomorfismo
canónico
\[
 D^i(S^{-1}A|R,S^{-1}M)
 \xrightarrow{\ \sim\ }S^{-1}D^i(A|R,M)
 \qquad(i\geq0).
\]

\item Sea $R$ noetheriano y sea $A$ un $R$-álgebra local
esencialmente de tipo finito. Denotemos por $\widehat A$ la
completación $\mathfrak m$-ádica de $A$, donde $\mathfrak m$ es el
ideal maximal de $A$. Entonces, para todo $A$-módulo $E$ de tipo
finito, hay un isomorfismo canónico
\[
 D^1(\widehat A|R,\widehat A\otimes_AE)
 \xrightarrow{\ \sim\ }
 \widehat A\otimes_AD^1(A|R,E).
\]
\end{enumerate}

\underline{Demostración.}

La parte (a) se deduce inmediatamente de
$D^1(A|R,M)=\operatorname{Exalcomm}(A|R,M)$.

Para (b), sea $P$ un álgebra de polinomios sobre $R$ en un número
finito de variables, y sea exacta
\[
 0\longrightarrow I\longrightarrow P\longrightarrow A
 \longrightarrow0.
\]
Después de localizar, podemos suponer que $I$ está generado por una
sucesión $P$-regular. Entonces $H_1(K_I)=0$, y por tanto
\[
 D^2(A|R,M)
 =\operatorname{Coker}\bigl(
 H^1(K_I,M)\longrightarrow
 \operatorname{Hom}_A(H_1(K_I),M)\bigr)=0.
\]
Recíprocamente, supongamos que $D^2(A|R,M)=0$ para todo $A$-módulo
$M$. Elijamos una sucesión exacta
\[
 0\longrightarrow K\longrightarrow F\longrightarrow I
 \longrightarrow0,
\]
donde $F$ es un $P$-módulo libre de tipo finito. Si
\[
 K_0=\operatorname{Im}\bigl(\textstyle\bigwedge^2F
 \longrightarrow\bigwedge^1F\bigr),
\]
entonces $H_1(K_I)=K/K_0$. Como
$D^2(A|R,H_1(K_I))=0$, el morfismo
\[
 H^1(K_I,H_1(K_I))
 \longrightarrow
 \operatorname{Hom}_A(H_1(K_I),H_1(K_I))
\]
es sobreyectivo. En consecuencia, la sucesión exacta
\[
 0\longrightarrow K/K_0\longrightarrow F/IF
 \longrightarrow I/I^2\longrightarrow0
\]
se escinde como sucesión de $A$-módulos. En particular, $I/I^2$ es
proyectivo. Después de localizar $A$ una vez más, podemos suponer que
$I/I^2$ es libre y que $F/IF\longrightarrow I/I^2$ es un
isomorfismo. La sucesión exacta escindida da entonces
$H_1(K_I)=K/K_0=0$. Así, $I$ está generado por una sucesión
$P$-regular, de modo que $A$ es una intersección completa relativa
sobre $R$.

La parte (c) se deduce inmediatamente de 3.12(4) y (6).

%%%%%%%%%%  scan idx 101  |  folio 90  %%%%%%%%%%

Para (d), 3.12(6) da primero isomorfismos canónicos
\[
 D^i(A|R,\widehat A\otimes_AE)
 \xrightarrow{\ \sim\ }
 \widehat A\otimes_AD^i(A|R,E)
\]
para todo $i$. Resta, pues, demostrar que
\[
 D^1(\widehat A|R,\widehat A\otimes_AE)
 \longrightarrow D^1(A|R,\widehat A\otimes_AE)
\]
es un isomorfismo. Utilizamos la descripción de $D^1$ en términos de
extensiones de álgebras.

Sea
\[
 0\longrightarrow\widehat A\otimes_AE\longrightarrow A'
 \xrightarrow{\ \varphi\ }\widehat A\longrightarrow0
\]
una extensión de $R$-álgebras cuya restricción a $A$ se escinde. Como
$\widehat A\otimes_AE$ es un $\widehat A$-módulo de tipo finito,
$A'$ es un $R$-álgebra local completa para su topología
$\mathfrak m'$-ádica, donde $\mathfrak m'$ es el ideal maximal de
$A'$. Si $\psi:A\longrightarrow A'$ es un morfismo de $R$-álgebras
que levanta el morfismo canónico $A\longrightarrow\widehat A$,
entonces $\psi$ se extiende de manera única a un morfismo
$\widehat\psi:\widehat A\longrightarrow A'$ que satisface
$\varphi\circ\widehat\psi=\operatorname{id}_{\widehat A}$.
Esto demuestra que
\[
 D^1(\widehat A|R,\widehat A\otimes_AE)
 \longrightarrow D^1(A|R,\widehat A\otimes_AE)
\]
es inyectivo.

%%%%%%%%%%  scan idx 102  |  folio 91  %%%%%%%%%%

Sea ahora
\[
 0\longrightarrow\widehat A\otimes_AE\longrightarrow B
 \xrightarrow{\ \varphi\ }A\longrightarrow0
\]
una extensión de $R$-álgebras, y pongamos
$\widehat E=\widehat A\otimes_AE$. Afirmamos primero que la topología
sobre $\widehat E$ inducida por la topología $\mathfrak m_B$-ádica de
$B$ coincide con su topología de $\widehat A$-módulo. En efecto, como
$A$ es esencialmente de tipo finito sobre $R$, existe una subálgebra
$B_0\subset B$ tal que $B_0$ es un $R$-álgebra local esencialmente de
tipo finito y $\varphi(B_0)=A$. Pongamos $E_0=\widehat E\cap B_0$, y
denotemos por $\mathfrak m_B$, $\mathfrak m_0$ y $\mathfrak m$ los
ideales maximales de $B$, $B_0$ y $A$, respectivamente. Entonces
\[
 \mathfrak m_B=\mathfrak m_0+\widehat E,
 \qquad
 \mathfrak m_B^{n+1}=\mathfrak m_0^{n+1}
   +\mathfrak m^n\widehat E
\]
para todo $n$. Como $B_0$ es noetheriano, el teorema de Artin--Rees
da un entero $r>0$ tal que, para todo $n\geq r$,
\begin{align*}
 \widehat E\cap\mathfrak m_B^{n+1}
 &=\mathfrak m^n\widehat E
   +\widehat E\cap\mathfrak m_0^{n+1}\\
 &=\mathfrak m^n\widehat E
   +E_0\cap\mathfrak m_0^{n+1}\\
 &=\mathfrak m^n\widehat E
   +\mathfrak m^{n+1-r}(E_0\cap\mathfrak m_0^r)
 \subset\mathfrak m^{n+1-r}\widehat E.
\end{align*}
Esto demuestra la afirmación sobre las dos topologías. En particular,
$\widehat E$ es completo para la topología inducida por la topología
$\mathfrak m_B$-ádica de $B$, y por tanto la completación da un
diagrama conmutativo exacto
\[
\begin{tikzcd}[column sep=large]
0 \arrow[r] & \widehat E \arrow[r] \arrow[d,equal]
& B \arrow[r] \arrow[d]
& A \arrow[r] \arrow[d] & 0 \\
0 \arrow[r] & \widehat E \arrow[r]
& \widehat B \arrow[r]
& \widehat A \arrow[r] & 0.
\end{tikzcd}
\]
Aquí $\widehat B$ denota la completación $\mathfrak m_B$-ádica de
$B$. Así,
\[
 D^1(\widehat A|R,\widehat E)
 \longrightarrow D^1(A|R,\widehat E)
\]
es sobreyectivo, y la demostración queda completa.
\hfill$\square$
